\documentclass[11pt]{article}
\usepackage[margin=1in]{geometry}
\usepackage{amsmath,amssymb,booktabs,url,xcolor}
\usepackage[hidelinks]{hyperref}
\usepackage{graphicx}
\newcommand{\pending}[1]{\textcolor{red}{[PENDING: #1]}}

\title{Staged Upgrades of Two Learned Populations: Keeping a Learned Fleet's Benefit When the Signals It Relies On Are Upgraded\\[4pt]
\large WORKING DRAFT (04/10/2026) --- structure written before results; every red item is a missing comparison}
\date{}

\begin{document}
\maketitle

\begin{abstract}
\pending{Destination unchanged: a method that is useful in the two-learned-population setting. Status: M1 met on grid4 (not replicated on
spider: the learned search fleet did not beat the rule there, and the rule's time-loss non-inferiority was not established on fresh traffic);
the M2 compatibility loss is replicated on spider with the rule fleet; the two light-side candidates confirmed so far (WC5, RPC) were worse
than the no-SPaT fallback; the retrained-fleet reference improves fuel but did not establish agent time-loss non-inferiority. No abstract
and no claim of a winning method until a candidate passes every gate against same-rights baselines.}
Setting sentence (allowed now): learned traffic signals and learned connected-vehicle fleets are deployed and upgraded on different
schedules; a fleet trained against one signal controller may lose its benefit when the controller is upgraded, and the operator can
retrain neither population at will.
\end{abstract}

\section{Introduction}
Problem; why staged upgrades; what is new must be stated against predictable-cycle RL signals (DHCP, arXiv 2509.03118),
signal--vehicle cooperative RL (CoTV; SVCC, TR-C 2026), robust eco-driving under uncertain SPaT (arXiv 1802.07192, 2208.13709) and
adaptive-SPaT provision (USDOT report). \pending{contribution list after M3; candidate: receiver-aware commitment.}

\section{Setting}
SUMO grid4 (16 signals, 1\,800 veh/h, 50\% fleet share, membership fixed per vehicle id); lights FT / ACT / learned (RL-H);
fleets: SUMO Krauss at imperfection $\sigma$ (non-learned family), ACC model, learned fleets (PPO, continuous acceleration).
Determinism and provenance: PYTHONHASHSEED = 0, bit-identical always-SUMO across jobs (C-27).

\section{Milestone M1: a learned fleet worth deploying}
\paragraph{Criterion (registered).} Agent fuel superiority vs always-SUMO with agent and human time-loss margins (+0.5 s/veh),
teleports and unfinished trips; strong baseline $\sigma^*$ chosen on SELECT; confirmation on four new policy seeds and fresh traffic.
\paragraph{Evidence so far.} No earlier recipe was confirmed (C-21 \dots C-31, corrected C-37); one development checkpoint (R5h\_t0)
passed on fresh traffic; the non-learned $\sigma$-family is a strong baseline ($\sigma=0.3$: $-2.5$ s/veh, $-2.2\%$ fuel).
\paragraph{Table 1 (fleet benefit; C-54, C-55).} Fresh traffic (20 seeds, generated after registration), FT, 50\% agents; learned
fleet = GLOSA-LD$\circ\sigma^*$, four policy parameters fitted by derivative-free search on training traffic (4 search seeds pooled).
Differences of the learned fleet, agent fuel / agent time loss / human time loss (95\% CI):
\begin{center}\small\begin{tabular}{lccc}\toprule
comparator & fuel (\%) & agent TL (s/veh) & human TL (s/veh)\\\midrule
always-SUMO & $-9.68$ [$-10.24$, $-9.00$] & $-3.48$ [$-3.77$, $-3.23$] & $-0.49$ [$-0.57$, $-0.41$]\\
$\sigma^*=0.3$ & $-7.58$ [$-8.14$, $-6.89$] & $-0.78$ [$-1.03$, $-0.58$] & $-0.36$ [$-0.45$, $-0.26$]\\
GLOSA (heuristic) & $-2.16$ [$-2.92$, $-1.35$] & $-3.03$ [$-3.88$, $-2.27$] & $-0.47$ [$-0.95$, $-0.07$]\\
GLOSA$\circ\sigma^*$ (no lane release) & $+0.16$ [$-0.57$, $+0.99$] & $-0.77$ [$-1.46$, $-0.16$] & $-0.30$ [$-0.65$, $+0.00$]\\
GLOSA-L$\circ\sigma^*$ (rule) & $-2.80$ [$-3.37$, $-2.06$] & $-0.18$ [$-0.41$, $+0.02$] & $-0.01$ [$-0.10$, $+0.08$]\\\bottomrule
\end{tabular}\end{center}
No teleports or unfinished trips for the learned fleet (GLOSA: 0.50, GLOSA$\circ\sigma^*$: 0.40 teleports per run). The fitted
parameters mainly learn where to release control (the last $\approx$15\,m before a needed lane change); the fuel gain over the rule is the
fuel the rule gives up by releasing early. Not RL; not compared with published eco-controllers; FT light only.
\paragraph{An RL fleet on the same panel (C-57).} A PPO speed residual on the rule (acting only where the rule commands; dense per-agent
cost; intervention-indexed discount, C-51), four new seeds pooled, against the rule: fuel $-1.70\%$ [$-2.33$, $-1.09$], agent time loss
$-0.04$ [$-0.29$, $+0.19$], human $+0.02$ [$-0.08$, $+0.12$]; against always-SUMO $-8.66\%$ / $-3.33$ / $-0.47$; against $\sigma^*$
$-6.54\%$ / $-0.64$ / $-0.33$; no teleports. It does not recover the fuel of GLOSA$\circ\sigma^*$ ($+1.29\%$ against it), which the fitted
release distance does.
\paragraph{Table 2 (SPaT ablation).} For the fitted GLOSA class, removing SPaT means the policy never commands and reduces to the
$\sigma^*$ fleet; the comparison against $\sigma^*$ above is therefore the class-level SPaT ablation ($-7.6\%$ fuel). It does not separate
information from representation. The PPO residual (R10, confirmed in M1, C-57) reduces to the $\sigma^*$ fleet exactly when no SPaT is
published (bit-identical rows, M2 check KM2-2), so its benefit against $\sigma^*$ is likewise the class-level SPaT ablation.
\paragraph{Diagnostics (appendix; C-39 \dots C-42).} D1: the ported phase reward (B) prefers physically better behaviour for large
differences but does not encode the time-loss constraints (it ranks an infeasible fuel-for-time policy far above always-SUMO); PPO from
scratch did not reach always-SUMO even on its training traffic. D2: a hand-written in-class controller (GLOSA target, no release) gives
$-6.7\%$ fuel at equal agent time but harms human cars; a behaviour clone of the heuristic GLOSA trades fuel for time in closed loop.
An evaluation bug (exit release missing in deterministic evaluation) was found by an identity check and fixed (C-41).
A residual on the in-class controller with a crawl-proof group cost (v9) failed its pilot gate; its confirmation was not evaluated (C-47).
Derivative-free policy search over a three-parameter GLOSA family on the deployment objective (fuel with time-loss penalties) returned
the untuned heuristic itself for the pilot search seed; the heuristic fails the $\sigma^*$ comparison on agent time loss (+1.8\,s/veh,
SELECT) and the human-time-loss margin on the width of its 5-seed bound (C-43).
\paragraph{An actuation pitfall shared by every commanded fleet (C-45).} With SUMO's \texttt{setSpeed}, every commanded fleet we
trained or wrote (GLOSA, an in-class controller, PPO, behaviour clones, earlier learned fleets) created teleports (2--17 per run), while
SUMO drivers at any $\sigma$ created none. Each inspected event is a lane-change deadlock: an agent car stopped at the end of a lane that does
not continue on its route. Releasing control whenever the car's current lane does not continue (a code predicate on SUMO's best-lanes
offset) removed these deadlocks on every panel tested (C-46). This is trace and intervention evidence in this SUMO version; without a
same-checkpoint factorial over the two interfaces we do not attribute the earlier learned fleets' time-loss failures to it in full (C-50).
\paragraph{A strong same-information rule (C-46).} GLOSA on the $\sigma^*$ base driver with lane-aware release (command only on a lane
that continues) is beneficial on fresh traffic (20 seeds): against $\sigma^*$, agent fuel $-5.0\%$ and agent / human time loss $-0.7$ /
$-0.4$\,s/veh; against always-SUMO, $-7.1\%$ and $-3.3$ / $-0.5$; no teleports. It is a fixed rule, not learned, not compared with
published eco-controllers, and without a SPaT ablation of its own it does not show that SPaT is necessary; it is the baseline a learned
fleet must beat on the same base and interface.
The learned four-parameter policy (R9) against this rule on fresh traffic is in Table~1 (C-54).

\section{Milestone M2: is there a compatibility problem?}
\paragraph{Design (target; partly run: M2a and M2b below; the adapted-fleet column is run as the retrained reference in M3; a learned $L_0$ has NOT been run).} Freeze the confirmed fleet $F_0$; lights $L_0$ (old) and $L_1$ (upgraded, learned), both beneficial in their
reference condition; matrix $L_0/F_0$, $L_1/F_0$, $L_0/F_1$, $L_1/F_1$ with absolute effects and difference-in-differences; SPaT
causal and correctly labelled (plan / guarantee / forecast with uncertainty).
\paragraph{M2a: SPaT service withdrawn or misused (C-59, C-60; fresh traffic, 20 seeds; four light seeds PAIRED with four fleet
seeds, not a factorial).} Scope of THIS panel: the old light is fixed-time (not a learned L0); it has no forecast arm (run later as M2b)
and no adapted-fleet column (run later as the retrained reference in M3). The shares of SPaT calls marked valid are calls, not measured forecast errors. The learned light is a large
upgrade over the fixed-time program with SUMO cars (city time loss $-33.7$\,s/veh, city fuel $-19.5\%$). Both confirmed fleets lose their
own benefit. If the learned light correctly publishes no plan, each fleet falls back to its base driver and becomes identical to
$\sigma^*$: the benefit drops from $-7.5\%$ (fitted fleet) and $-6.5\%$ (PPO fleet) agent fuel to exactly 0, with no harm. If the fleet
keeps deriving SPaT from the static program, which the learned light no longer follows (30--43\% of its SPaT calls are marked valid
outside the function's guaranteed scope; these are calls, not measured forecast errors), it is harmed against $\sigma^*$ under the same light: $+7.4\%$ / $+6.5\%$ fuel and $+1.7$ / $+1.6$\,s/veh agent time loss. In
absolute terms the agents still gain from the upgrade; what is lost is the fleet's benefit against its counterfactual drivers.
\paragraph{M2b: an empirical forecast from the learned light (C-61, C-71, C-72).} A conditional-quantile forecast of the learned light's
green durations does not return the benefit to the frozen fleet. Two facts explain why this is hard: the original forecast was 1\,s early in
the green branch on 100\% of points (a timing-convention error, fixed in a corrected provider; one yellow branch remains 1\,s early, C-84),
and the frozen fleet is sensitive to timing errors, with asymmetric fuel and time effects in an open-loop oracle-shift diagnostic (exact
timing $-6.6\%$ fuel; 2\,s early $-2.4\%$; 2\,s late $-4.2\%$ fuel but agent time loss $+1.5$\,s/veh; 8\,s late $+9.3\%$ and $+11.6$\,s/veh;
C-72). The diagnostic uses another fleet and schedule than the retrained fleet, so it warns of the trade-off but does not predict it.
\paragraph{Second network (C-81).} On an independent own-built network (SUMO \texttt{netgenerate} spider, 18 signals, six with three greens),
upgrading the static programs to a learned all-green controller cuts city time loss by 58.9\,s/veh. The frozen rule fleet (GLOSA-L$\circ\sigma^*$,
$-9.2\%$ agent fuel under the static programs) again loses its benefit: labelled $\to$ exactly $\sigma^*$ (no harm); naive $\to$ harmed
($+7.2\%$ fuel, agent TL $+2.7$, human TL $+0.6$, upper $+0.8$).
\begin{figure}[t]\centering\includegraphics[width=\linewidth]{figs_two_pop/fig1_m2_crossplay.pdf}
\caption{The frozen fleet's benefit against $\sigma^*$ under the same light, before (FT) and after the learned-light upgrade. Labelled bars are
exactly 0 (rows identical to $\sigma^*$). Crossed 95\% intervals; grid4 fresh 470--489 (four paired light/fleet seeds), spider fresh
1000--1019 (two light seeds). Script \texttt{scripts/CITY\_fig\_m2\_crossplay.py}.}\end{figure}
\paragraph{Known so far (not the target experiment).} Weak learned fleets regressed under light upgrades and a car-side reaction repair
helped (C-14, C-16); non-learned fleet upgrades ($\sigma$, ACC) were not harmed by three learned-light families (C-30, C-32 \dots C-34).
If $F_0$ shows no regression, the hazard hypothesis is rejected in this setting.

\section{Milestone M3: light-side certainty vs the no-SPaT fallback (candidates tested so far)}
\paragraph{Question.} Can the upgraded learned light return its benefit to the FROZEN fleet by making its timing certain or learnable, at
less light cost than the strongest commitment baseline? Endpoint: agent fuel of the frozen fleet against R, the learned light publishing no
SPaT (the fleet then drives as $\sigma^*$); time-loss constraints reported.
\paragraph{Candidates screened on DEV (grid4, light seeds 0--1, two dev fleets).} Reward-shaped lights
that are paid for the fleet's forecast accuracy (E family, 10 arms; C-68), fixed-window commitment W-c ($W\in\{5,10,15\}$), ETA-triggered
commitment B, a predictable cycle, a policy-aware forecast, and receiver-triggered precision commitment RPC (C-79; 0 violations in 247\,498
commitments). None passed its DEV gate (agent fuel superiority over R within the time-loss limits). Not running their confirmations was
a post-DEV decision (C-69, C-79); the negative confirmation below was then registered separately.
\paragraph{Confirmed negative (M3-NC, C-86).} On fresh traffic, five new learned-light seeds and four frozen fleets (crossed, 20 cells):
fixed-window commitment WC5 raises agent fuel by $+2.98\%$ $[+1.64, +4.07]$ and RPC by $+9.26\%$ $[+8.51, +9.95]$; every one of the 20
cells has a positive point estimate (the cells share checkpoints; they are a counterexample search, not independent replications; the
interval comes from the crossed bootstrap). A shadow fleet (the light computes its messages exactly as in deployment, while the fleet drives
as $\sigma^*$) gives a deployment decomposition on a common denominator: running the mechanism costs $+5.55\%$ / $+10.59\%$ agent fuel
($+5.6$ / $+7.2$\,s/veh city time loss), and activating the fleet under the mechanism then reduces fuel by $2.57$ / $1.33$ points relative to its
matched shadow. The second term is the fleet's closed-loop contribution under that mechanism (the light also reacts to the fleet), not an
isolated effect of the messages. Replacing WC5's uncorrected forecast fallback by the corrected
one reduces the loss to $+0.43\%$ $[-1.21, +2.02]$ (descriptive), still with $+5.5$\,s/veh city time loss.
\paragraph{Second network (C-82).} Fixed-window commitment on the spider network (the stage-1 lights and the stage-3 panel, rule fleet):
$+6.68\%$ agent fuel, mechanism cost $+6.61\%$, fleet contribution $+0.06\%$ $[-0.27, +0.41]$. No fleet contribution is detected there, so the
grid4 reading (a fleet contribution smaller than the mechanism cost) does not transfer; on spider the whole loss is the mechanism cost.
\paragraph{Reference with update rights: retraining the fleet (C-78, C-85).} A fleet retrained from zero under the upgraded light, with the
forecast provider, saves $-2.61\%$ $[-3.06, -2.16]$ agent fuel on five new lights $\times$ four fleets $\times$ fresh traffic, with city and
human time loss within the margin. Agent time loss rises by $+0.34$\,s/veh (two-sided interval $[+0.08, +0.59]$, one-sided upper $+0.55$):
non-inferiority at the registered $+0.5$\,s/veh margin was not established, which is not evidence that the increase exceeds the margin
(C-85). Not confirmed. This differs from WC5 / RPC, which are significantly worse than the fallback.
The same rule without learning is $+3.18\%$.
\begin{figure}[t]\centering\includegraphics[width=\linewidth]{figs_two_pop/fig2_m3_frontier.pdf}
\caption{Cost--contribution comparison (the points tested; not an optimal frontier). A: both components on the denominator R in every cell,
so they add exactly to the endpoint (net labels; crossed 95\% intervals). WC5 and RPC: grid4, fresh 510--529, five new lights, four frozen
fleets. WCK5: spider, stage-1 lights, stage-3 panel, rule fleet. F1: grid4 retrained fleet (more update rights, no light change, training cost
not shown), fresh 530--549; it improves fuel but did not establish agent time-loss non-inferiority. The shaded region is lower fuel only, not
the deployment gate. B: DEV screens in the DEV reports' units (descriptive). Script \texttt{scripts/CITY\_fig\_m3\_frontier.py} (v2).}\end{figure}
\paragraph{Reading.} Only the candidates tested are closed. For WC5 and RPC on grid4 the commitments were honoured (0 violations), and
the mechanism cost exceeded the fleet's closed-loop contribution; on spider no fleet contribution was detected. Retraining the fleet improves
fuel but did not establish agent time-loss non-inferiority. This is not a claim that no light-side mechanism can work.

\section{Milestone M4: a second setting}
Spider network (own-built, \texttt{netgenerate}): a learned all-green light beats the static programs by 57\,s/veh (two seeds, C-76); the
rule fleet saves $-9.2\%$ fuel against $\sigma^*$ (C-77), with agent / human time-loss non-inferiority not established on fresh traffic
(upper $+0.55$ / $+0.52$ against the $+0.5$ margin, C-81);
the four-parameter search fleet (R9) does NOT beat the rule there (C-80); the M2 loss and harm replicate with the rule fleet (C-81); fixed-window commitment is again worse than the
no-SPaT fallback, with a different decomposition (no fleet contribution detected, C-82). The retrained-fleet check was not run there, because its registered condition (F1-A1 confirmed) failed. The
earlier RESCO cologne8 attempt failed at the light baseline and is kept as a negative (C-35, C-36).

\section{Limitations}
Simulation only; one fuel model (SUMO HBEFA); FT used for the fleet pilot; the reward-B port: fuel enters only through the $\delta$
threshold in this unit choice (calibration, Amendment R3).
One learned-light family per network; $W=5$ only in confirmation. The convention-corrected provider still has a 1\,s early yellow branch
(C-84), identical in training and evaluation of the retrained fleet. The spider fleet is a rule, not a learned fleet.

\appendix
\section{Validation and pitfalls (QA)}
Reward traces, sentinel readings, residual ratchet, horizon semantics, dominance labels (C-17, C-22, C-25, C-27, C-37, C-38).

\end{document}
